\section{Linear machines: the linear crux as rigidity}\label{sec:linear}

For machines whose data path is linear over a field, Theorem~D has a linear analogue, Theorem~$\Dlin$ (\cref{thm:Dlin}):
\emph{linear depth-TCP} together with a \emph{linear crux} excludes oblivious linear-time transposers. Linear depth-TCP weakens
the structural hypothesis, replacing depth by \emph{linear} depth, in which long paths may cancel; the linear crux is required
for a larger class of systems, the ball systems of the schedules, and it can be stated exactly as relative rigidity of the
transposition matrix with respect to the dependency pattern of the schedule. Relative rigidity alone excludes no machine.

\subsection{Linear data paths and path sums}\label{sec:lin-setting}

Let $F$ be a field; input bits are read as $0,1\in F$. Let $M$ be oblivious at length~$L$, with dependency DAG $H_M$ and step
values $v_t(w)$ as in \cref{sec:prelim-tcp}, and write $\mathrm{out}(t)$ for the set of defined out-neighbours $t-1$,
$\prev_1(t),\dots,\prev_k(t)$ of~$t$ in~$H_M$.

\begin{definition}[Linear data path]\label{def:linear-path}
\partformal{Rigidity.lean}{LinNet, LinNet.val}{the machine, its step values and their encodings; in Lean the data path is
given as a network, and the output steps and the functionals $\varphi_\ell$ are arguments of \lean{LinNet.rigidity}.}
$M$ has an \emph{$F$-linear data path} at length~$L$ if there are maps $\mathrm{enc}_t$ from the possible values of step~$t$ to
$F^m$, vectors $c_t,a_t\in F^m$ and matrices $A_{t,s}\in F^{m\times m}$, none depending on the input, such that for every
$w\in\{0,1\}^L$
\[
  \mathrm{enc}_t(v_t(w))=c_t+\sum_{s\in\mathrm{out}(t)}A_{t,s}\,\mathrm{enc}_s(v_s(w))+a_t\,w_{x_0(t)},
\]
where the last term is present only if the input head reads a position $x_0(t)\in[0,L)$, and such that output bit~$\ell$,
appended at step~$t_\ell$, equals $\varphi_\ell(\mathrm{enc}_{t_\ell}(v_{t_\ell}(w)))$ for a fixed affine functional
$\varphi_\ell$ with linear part $\varphi_\ell^{\mathrm{lin}}$. Such a \emph{representation} need not be unique.
\end{definition}

Write $v_t$ for $\mathrm{enc}_t(v_t)$. For a genie~$\Gamma$, expanding the recursion along paths and stopping at steps
of~$\Gamma$ gives
\[
  v_t=c'_t+\sum_{s\in\Gamma}T_{t,s}v_s+\sum_jC_{t,j}w_j .
\]
Here, for a path $\pi=(t=t_0\to\dots\to t_n)$ of $H_M$ with weight $W(\pi)=A_{t_0,t_1}\cdots A_{t_{n-1},t_n}$, $T_{t,s}$ is
the sum of $W(\pi)$ over the paths from~$t$ to $s\in\Gamma$ with $n\ge1$ whose interior avoids~$\Gamma$; $C_{t,j}$ is the sum
of $W(\pi)a_{t_n}$ over the $\Gamma$-free paths whose last step reads position~$j$; and $c'_t$ is the sum of $W(\pi)c_{t_n}$
over all $\Gamma$-free paths from~$t$. The sums are finite because $H_M$ is acyclic. In matrix form they are the blocks of
$(I-A_\Gamma)^{-1}(A-A_\Gamma)$ and $(I-A_\Gamma)^{-1}B$, where $A_\Gamma$ is the block matrix~$A$ with the columns of~$\Gamma$
set to zero and $B$ is the input matrix: the transfer matrices of the linear network cut at~$\Gamma$, in the sense of Koetter and
Médard~\cite{KoetterMedard2003}.

\begin{definition}[Linear depth]\label{def:linear-depth}
\partformal{Rigidity.lean}{LinNet.GFree, LinNet.C, LinNet.reach}{the condition of linear depth itself, which in Lean is a
hypothesis of \lean{LinNet.rigidity}.}
Let $I^D_\Gamma(\ell)$ be the set of input positions read at steps reachable from~$t_\ell$ by $\Gamma$-free paths of length
at most~$D$, and $B^D(\ell)\supseteq I^D_\Gamma(\ell)$ the set of positions read at steps reachable from~$t_\ell$ by directed paths of
length at most~$D$ in~$H_M$, so that $B^D(\ell)=I^D_\emptyset(\ell)$. Output step~$t_\ell$ has \emph{linear depth at most~$D$ for~$\Gamma$} if
$\varphi_\ell^{\mathrm{lin}}C_{t_\ell,j}=0$ for every $j\notin I^D_\Gamma(\ell)$.
\end{definition}

For such~$j$ every $\Gamma$-free path to a read of~$j$ is longer than~$D$, so the condition says that the long contributions
to inputs outside $I^D_\Gamma(\ell)$ cancel; output bit~$\ell$ is then an affine function of the genie values and of the
bits~$w_j$, $j\in I^D_\Gamma(\ell)$. Structural depth at most~$D$ implies linear depth at most~$D$ at every output step.
Linear depth is not monotone in~$\Gamma$: adding a step to~$\Gamma$ can break a cancellation between a path through it and a
path that avoids it.

\subsection{Rigidity relative to the schedule}\label{sec:lin-rigid}

Let $P_{\tpos}$ be the $L\times L$ matrix whose row~$\ell$ is the unit vector $e_{\tpos(\ell)}$. The \emph{minrank} of a
system $(I(\ell))_\ell$ over~$F$ is the least rank of a matrix~$R$ such that row~$\ell$ of $P_{\tpos}-R$ is supported in
$I(\ell)$ for every~$\ell$; it is the least length of an $F$-linear index code for the system, decodable on all of~$F^L$ with
linear decoders~\cite{Haemers1979,BBJK2011}.

\begin{lemma}[Per-machine rigidity]\label{lem:rigidity}
\partformal{Rigidity.lean}{LinNet, LinNet.val, LinNet.C, LinNet.reach, LinNet.reach_subset, LinNet.Pmat, LinNet.minrank, LinNet.rigidity, LinNet.minrank_le}{the
translation from the machine to a linear network: the machine model, the encodings $\mathrm{enc}_t$ and the step
values~$v_t$.}
Suppose that $M$ has an $F$-linear data path at length~$L$ (with some representation), transposes all $d\times d$ matrices,
and has a genie~$\Gamma$ under which every output step has linear depth at most~$D$. Then
\[
  P_{\tpos}=R+S,\qquad\rank R\le\beta\card\Gamma,\qquad\supp(\text{row }\ell\text{ of }S)\subseteq I^D_\Gamma(\ell)
  \subseteq B^D(\ell);
\]
equivalently, the system $(I^D_\Gamma(\ell))_\ell$ has minrank at most $\beta\card\Gamma$ over~$F$.
\end{lemma}

\begin{outline}
The Lean statement concerns a linear network, \lean{LinNet}, with out-neighbours at earlier steps, the vectors $c_t$ and
$a_t$, the matrices $A_{t,s}$, and the input position read by each step, if any. Its formal values $V_t(w)$, $w\in F^L$, are
defined by the recursion of the data path (\lean{LinNet.val}), and the hypotheses are stated for them: on binary inputs the
affine output functionals give the transposed bits, each step of~$\Gamma$ takes at most $2^\beta$ values, and every output
step has linear depth at most~$D$, with the path sums \lean{LinNet.C} over the $\Gamma$-free paths (\lean{LinNet.GFree}) and
the sets $I^D_\Gamma(\ell)$ given by \lean{LinNet.reach}; for $\Gamma=\emptyset$ these are the ball sets, and
$I^D_\Gamma(\ell)\subseteq B^D(\ell)$ (\lean{LinNet.reach_subset}).

The values split into a constant and a linear part, $V_t(w)=V_t(0)+V^{\mathrm{lin}}_t(w)$ (\lean{LinNet.val_eq_add}), and the
linear part has the path-sum expansion
$V^{\mathrm{lin}}_t(w)=\sum_{s\in\Gamma}T_{t,s}V^{\mathrm{lin}}_s(w)+\sum_jC_{t,j}w_j$ identically in $w\in F^L$, being the
unrolled recursion (\lean{LinNet.C_eq}, \lean{LinNet.vlin_expand}); in Lean $T_{t,s}$ is defined by the recursion that it
satisfies (\lean{LinNet.Tm}). Let $E(w)=(V^{\mathrm{lin}}_s(w))_{s\in\Gamma}$, the linear parts of the revealed values, let
$R$ be the matrix of the linear map $w\mapsto\bigl(\sum_{s\in\Gamma}\varphi_\ell^{\mathrm{lin}}T_{t_\ell,s}E_s(w)\bigr)_\ell$,
and let $S_{\ell j}=\varphi_\ell^{\mathrm{lin}}C_{t_\ell,j}$. Transposition at the inputs $0$ and~$e_j$ gives
$\varphi_\ell^{\mathrm{lin}}V^{\mathrm{lin}}_{t_\ell}(e_j)=(e_j)_{\tpos(\ell)}$, and the expansion turns this into
$(P_{\tpos})_{\ell j}=R_{\ell j}+S_{\ell j}$. By linear depth, $S_{\ell j}=0$ for $j\notin I^D_\Gamma(\ell)$. On $\{0,1\}^L$,
$E$ takes at most $2^{\beta\card\Gamma}$ values, since each $V_s$, $s\in\Gamma$, takes at most $2^\beta$ values there. A
linear map of rank~$\rho$ takes at least $2^\rho$ values on $\{0,1\}^L$: it is injective on the binary vectors supported on
$\rho$ coordinates whose columns are independent, because differences of such vectors have coefficients in $\{-1,0,1\}$
(\lean{LinNet.two_pow_finrank_le}). So $\rank R\le\rank E\le\beta\card\Gamma$ (\lean{LinNet.rigidity}); in other words, $E$
is an $F$-linear index code of rank at most $\beta\card\Gamma$ for $(I^D_\Gamma(\ell))_\ell$, decodable with linear decoders.
The minrank form is \lean{LinNet.minrank_le}, with the minrank as the least such rank (\lean{LinNet.minrank}).
\end{outline}

\begin{proof}[Proof of the translation (not formalized)]
The dependency DAG of~$M$, with out-neighbours $\mathrm{out}(t)$, the data $c_t,A_{t,s},a_t$ of the representation and the
input positions read define a linear network. By induction on~$t$, its formal values satisfy $V_t(w)=\mathrm{enc}_t(v_t(w))$
for $w\in\{0,1\}^L$, since both satisfy the recursion. As $M$ transposes, it appends exactly $L$ bits, and output bit~$\ell$
equals $\varphi_\ell(V_{t_\ell}(w))$ on binary inputs; each $v_s$ takes at most $2^\beta$ values; and linear depth is a
statement about the path sums of the representation.
\end{proof}

\begin{definition}[Linear depth-TCP and the linear crux]\label{def:lin-hyp}
\emph{Linear depth-TCP}$(k,C,\beta,\eps,D;F)$ says: for every machine~$M$ with $k$ work tapes and value width at most~$\beta$
there is~$d_0$ such that at every $d\ge d_0$ at which $M$ is oblivious at length~$L$, has an $F$-linear data path and halting
time at most $C(L+1)$, there are a representation of its data path and a genie~$\Gamma$ with $\beta\card\Gamma\le\eps L$ under
which every output step has linear depth at most~$D$. Let $\Sball_{k,C,\beta,D}$ be the class of the ball systems
$(B^D(\ell))_{\ell\in[0,L)}$, at the lengths $L=d^2b(d)$, of the oblivious machines with $k$ work tapes, value width at
most~$\beta$ and halting time at most $C(L+1)$ that append exactly $L$ bits, as for $\mathcal S_{k,C,\beta,D}$; no linear data
path is required. $\LinCrux_F(\Sball_{k,C,\beta,D},\rho)$ says: for every $\eta>0$ there is~$d_0$ such
that every system of $\Sball_{k,C,\beta,D}$ at a size $d\ge d_0$ has minrank at least $(\rho-\eta)L$ over~$F$.
\end{definition}

\begin{theorem}[Theorem $\Dlin$]\label{thm:Dlin}
\partformal{Rigidity.lean}{LinNet.dlin_step}{the quantification over machines and sizes, the two hypotheses as statements
about machines, and the translation of a machine into a linear network (\cref{lem:rigidity}); the Lean theorem is the
argument at one size, for a network.}
Let $k\ge1$, $D\ge0$ and $\beta\ge2$ be integers, $C>0$, $0<\eps<\rho$ and $F$ a field. If linear depth-TCP$(k,C,\beta,\eps,D;F)$
and $\LinCrux_F(\Sball_{k,C,\beta,D},\rho)$ hold, then for every machine~$M$ with $k$ work tapes and value width at
most~$\beta$ there is~$d_0$ such that for no $d\ge d_0$ is $M$ oblivious at length~$L$ with an $F$-linear data path, halting
time at most $C(L+1)$, and a transposer of all $d\times d$ matrices.
\end{theorem}

\begin{outline}[Outline of the formal proof at one size]
For a transposing network with a genie~$\Gamma$ of cost $\beta\card\Gamma\le\eps L$ under which every output step has
linear depth at most~$D$, \cref{lem:rigidity} bounds the minrank of the ball system by $\beta\card\Gamma\le\eps L$; with
$\eta=(\rho-\eps)/2$ and $L\ge1$ this contradicts a minrank of at least $(\rho-\eta)L$ (\lean{LinNet.dlin_step}).
\end{outline}

\begin{proof}[Proof of the reduction to one size (not formalized)]
Let $\eta=(\rho-\eps)/2$ and let $d_0$ be the larger of the thresholds of the two hypotheses. If $M$ were such a transposer at
some $d\ge d_0$, linear depth-TCP would give a representation and a genie~$\Gamma$ with $\beta\card\Gamma\le\eps L$ under which
every output step has linear depth at most~$D$, and $\LinCrux_F$ would give the ball system
$(B^D(\ell))_\ell\in\Sball_{k,C,\beta,D}$ a minrank of at least $(\rho-\eta)L$; the argument at one size, applied to the
network of the representation, gives a contradiction.
\end{proof}

This is the proof of Theorem~D with the cone lemma replaced by the path-sum expansion, which is exact for linear data paths.

\paragraph{Why it matters.} Read as a matrix statement, $\LinCrux_F(\Sball,\rho)$ says that $P_{\tpos}$ is $\rho L$-rigid
relative to the $D$-ball pattern of every linear-time schedule: it is not the sum of a matrix of rank below $\rho L$ and a
matrix supported on that pattern. So, for linear machines, linear depth-TCP together with a
linear crux excludes oblivious linear-time transposers, and the linear crux can be stated exactly as relative rigidity of the
transposition matrix; \cref{lem:rigidity} applies only to a machine that has a cheap genie under which every output has
linear depth at most~$D$. This is a relative form of Valiant's route from depth reduction and rigidity to lower bounds for
linear circuits~\cite{Valiant1977}. The route from small depth to index coding also goes back to Riis's guessing
games~\cite{Riis2007}, and it was used in the language of min-rank by Jukna and Schnitger~\cite{JuknaSchnitger2011}. For
static data structures, lower bounds imply rigidity~\cite{DGW2019}, and for systematic linear data structures the two are
equivalent~\cite{NRR2020}. Without the pattern the statement is false for permutation matrices: splitting the receivers into
groups of $s+1$, each knowing the bits demanded by the others in its group, gives rank $\lceil L/(s+1)\rceil$ with $s$
off-diagonal entries per row~\cite[Section~4.3]{FourthPaper2026}. Any proof must use the schedule's structure. On the construction
side, a linear oblivious transposer at size~$d$ either has no genie of cost $\eps L$ that makes every output linearly
$D$-shallow, or its formal revealed map is a linear index code of rank at most $\eps L$ for its own ball system.

\begin{remark}[Scope]\label{rem:Dlin-scope}
(a) The genie plays no role in the crux class: since $I^D_\Gamma(\ell)\subseteq I^D_\emptyset(\ell)=B^D(\ell)$ and the empty
genie is allowed, the linear crux for the systems $(I^D_\Gamma(\ell))_\ell$ over all genies is equivalent to the linear crux
for the ball systems; and $\Sball$ dominates the fourth paper's class $\mathcal S_{k,C,\beta,D}$ pointwise. Neither
Theorem~$\Dlin$ nor the linear form of Theorem~D implies the other: $\Dlin$ trades a weaker structural hypothesis for a larger
crux class. For home-based linear machines $\Dlin$ gives nothing new, since Corollary~B supplies structural genies.
(b) The trade-off between cost and depth survives: the machine that appends, at every step, the parity of all bits read so
far forces at least $\lfloor(L-1)/(D+1)\rfloor$ revealed steps for linear depth~$D$, as in the fourth paper's forced trade-off
\cite[Proposition~4.13]{FourthPaper2026}.
(c) Linear depth is robust to dummy traffic on separate coordinates: copies of an input bit carried along deep paths of a
dummy schedule on extra tapes and combined so that they cancel contribute $W(\pi_1)+W(\pi_2)=0$ to the path sums, whatever their
lengths, provided the genie reveals single coordinates (fields, as in~\cite[Remark~4.6]{FourthPaper2026}) and does not cut one
of the two paths; Lemma~\ref{lem:rigidity} and Theorem~$\Dlin$ hold verbatim for such coordinate genies, with cost equal to the
total bit width of the revealed coordinates. With whole-step genies the robustness fails, and it fails for dummy traffic that is
entangled with revealed coordinates, for instance a value masked by a pad and unmasked through a deep partner copy.
(d) Linear depth-TCP quantifies over all linear machines, and the theorem uses it only for the transposer itself. It is not a
statement about head schedules alone: it depends on the data path and on its representation.
\end{remark}
